\handout
{Collection of Problems}
{\textsc{Single Variable Calculus (SVC)}}
{\href{https://creativecommons.org/publicdomain/zero/1.0/}{CC0 1.0 Universal license}}
{Author: Jaehoon Song}{Release: \generatedOn}
{SVC: Collection of Problems}


% \begin{subanswer}[red][true]
% \begin{subanswer}[red][true][red][true]

\begin{enumerate}
  \item \textbf{Particular Solution Parameter} (10 points)\\
  For what value of $k$, if any, is
  \[
    y = e^{-2x} + k e^{4x}
  \]
  a solution to the differential equation
  \[
    y - \frac{y''}{4} = 5 e^{4x}?
  \]
  \begin{subanswer}
    % your answer here
  \end{subanswer}

  \item \textbf{Tangent Line Interpretation} (10 points)\\
  An equation for the line tangent to the graph of the differentiable function $f$ at $x = 3$ is
  \[
    y = 4x + 6.
  \]
  Which of the following statements \emph{must} be true?\\
  \begin{enumerate}[label=(\Roman*)]
    \item $f(0) = 6$
    \item $f(3) = 18$
    \item $f'(3) = 4$
  \end{enumerate}
  \begin{subanswer}
    % your answer here
  \end{subanswer}

  \item \textbf{General Solution via Separation of Variables} (10 points)\\
  What is the general solution to the differential equation
  \[
    \frac{dy}{dx} = \frac{x \cos\bigl(x^2\bigr)}{4y},
    \quad y > 0?
  \]
  \begin{subanswer}
    % your answer here
  \end{subanswer}

  \item \textbf{Derivative of an Inverse Tangent} (5 points)\\
  Compute
  \[
    \frac{d}{dx}\bigl(\tan^{-1}(3x)\bigr).
  \]
  \begin{subanswer}
    % your answer here
  \end{subanswer}

  \item \textbf{Squeeze Theorem Limit} (10 points)\\
  Let
  \[
    g(x) = \sin\Bigl(\tfrac{\pi}{2}(x+2)\Bigr) + 3,\quad
    h(x) = -\tfrac{1}{4}x^3 - \tfrac{3}{2}x^2 - \tfrac{9}{4}x + 3.
  \]
  If $f$ is a function satisfying
  \[
    g(x) \le f(x) \le h(x)
    \quad\text{for } -2 < x < 0,
  \]
  what is
  \[
    \lim_{x \to -1} f(x)?
  \]
  \begin{subanswer}
    % your answer here
  \end{subanswer}

  \item \textbf{Implicit Differentiation} (5 points)\\
  If
  \[
    y = \ln\bigl(3x + 4y\bigr),
  \]
  then compute
  \[
    \frac{dy}{dx}.
  \]
  \begin{subanswer}
    % your answer here
  \end{subanswer}

  \item \textbf{Riemann Sum to Integral} (10 points)\\
  Which of the following definite integrals are equal to
  \[
    \lim_{n \to \infty} \sum_{k=1}^n \sin\Bigl(-1 + \tfrac{5k}{n}\Bigr) \frac{5}{n}?
  \]
  \begin{enumerate}[label=(\Roman*)]
    \item $\displaystyle \int_{-1}^{4} \sin x \,dx$
    \item $\displaystyle \int_{0}^{5} \sin(-1+x) \,dx$
    \item $\displaystyle 5 \int_{0}^{1} \sin(-1 + 5x) \,dx$
  \end{enumerate}
  \begin{subanswer}
    % your answer here
  \end{subanswer}

  \item \textbf{Chain Rule Application} (10 points)\\
  If
  \[
    f(x) = \bigl(e^{3x} + \sin(2x)\bigr)^4,
  \]
  then compute
  \[
    f'(x).
  \]
  \begin{subanswer}
    % your answer here
  \end{subanswer}

  \item \textbf{Definite Integral Evaluation} (10 points)\\
  Evaluate
  \[
    \int_{3}^{5} \frac{2x^2 + x + 4}{x - 1} \,dx.
  \]
  \begin{subanswer}
    % your answer here
  \end{subanswer}

  \item \textbf{Product Rule Derivative} (5 points)\\
  Compute
  \[
    \frac{d}{dx}\bigl(\cos x \tan x\bigr).
  \]
  \begin{subanswer}
    % your answer here
  \end{subanswer}

  \item \textbf{Inverse Function Derivative} (10 points)\\
  For which of the following decreasing functions $f$ does
  \[
    \bigl(f^{-1}\bigr)'(10) = -\tfrac{1}{8}?
  \]
  \begin{enumerate}[label=(\Alph*)]
    \item $f(x) = -5x + 15$
    \item $f(x) = -2x^3 - 2x + 14$
    \item $f(x) = -x^5 - 4x + 15$
    \item $f(x) = e^{-2x} - x + 9$
  \end{enumerate}
  \begin{subanswer}
    % your answer here
  \end{subanswer}

  \item \textbf{Exponential Decay Model} (10 points)\\
  A dose of 400 milligrams of a drug is administered to a patient. The amount of the drug, $A(t)$, in the bloodstream satisfies
  \[
    \frac{dA}{dt} = k A,
  \]
  for constant $k$. Which of the following could be an expression for $A(t)$?\\
  \begin{enumerate}[label=(\Alph*)]
    \item $A(t) = 400 e^{-0.3t}$
    \item $A(t) = e^{-0.3t} + 399$
    \item $A(t) = -3t + 400$
  \end{enumerate}
  \begin{subanswer}
    % your answer here
  \end{subanswer}

  \item \textbf{Continuity Check} (5 points)\\
  Confirm that the function
  \[
    f(x) = \frac{x^2 - 5x + 4}{x^2 - 6x + 8}
  \]
  is continuous at $x = 4$.  
  \begin{subanswer}
    % your answer here
  \end{subanswer}

  \item \textbf{Instantaneous Rate at a Time} (5 points)\\
  The depth of water $t$ hours after midnight is
  \[
    D(t) = 10 + 4.9 \cos\Bigl(\tfrac{\pi}{6} t\Bigr).
  \]
  Describe the method to find the instantaneous rate of change (in ft/hr) at 6 A.M.
  \begin{subanswer}
    % your answer here
  \end{subanswer}

  \item \textbf{Derivative at a Point} (5 points)\\
  Let
  \[
    f(x) = 2x^3 - 4x^2 - 5.
  \]
  What is $f'(-1)$?
  \begin{subanswer}
    % your answer here
  \end{subanswer}

  \item \textbf{Indefinite Integral} (10 points)\\
  Compute
  \[
    \int \frac{12 x^2}{2x + 1} \,dx.
  \]
  \begin{subanswer}
    % your answer here
  \end{subanswer}

  \item \textbf{Limit Evaluation} (5 points)\\
  Evaluate
  \[
    \lim_{x \to 3} \frac{x - 3}{x^3 - 9x}.
  \]
  \begin{subanswer}
    % your answer here
  \end{subanswer}

  \item \textbf{Combination of Derivatives} (5 points)\\
  If $f'(1) = 2$ and $g'(1) = 6$, and
  \[
    h(x) = 5f(x) - 4g(x) + 3x^2 - 2,
  \]
  find $h'(1)$.
  \begin{subanswer}
    % your answer here
  \end{subanswer}

  \item \textbf{Differential Equation Solutions} (5 points)\\
  What are all solutions to
  \[
    \frac{dy}{dx} = \frac{1}{x \ln 2}?
  \]
  \begin{subanswer}
    % your answer here
  \end{subanswer}

  \item \textbf{Related Rates Relationship} (10 points)\\
  Given $U = \tfrac{k}{N}$ with constant workforce, describe the relationship between $\tfrac{dU}{dt}$ and $\tfrac{dN}{dt}$.
  \begin{subanswer}
    % your answer here
  \end{subanswer}

  \item \textbf{Points of Inflection} (10 points)\\
  Find the $x$-values where
  \[
    y = e^{-x} + 2x e^{-x} + x^2 e^{-x}
  \]
  has inflection points.
  \begin{subanswer}
    % your answer here
  \end{subanswer}

  \item \textbf{Vertical Tangents} (10 points)\\
  Describe the $y$-coordinates of points on the curve
  \[
    e^x = \sin y
  \]
  where the tangent line is vertical.
  \begin{subanswer}
    % your answer here
  \end{subanswer}










  %% continues on ... Q30
  \item \textbf{Inverse Function Derivative} (10 points)\\
  Let $f$ and $g$ be inverse functions that are differentiable for all $x$. If $f(-5)=7$ and $g^{\prime}(7)=3$, which of the following statements must be false?
  \begin{enumerate}[label=(\Roman*)]
    \item $f^{\prime}(3)=-\frac{1}{3}$
    \item $f^{\prime}(-5)=\frac{1}{3}$
    \item $f^{\prime}(7)=\frac{1}{3}$
  \end{enumerate}
  \begin{subanswer}
    % Student’s answer space
  \end{subanswer}

  \item \textbf{Marginal Cost Estimate} (10 points)\\
  The function $C$ gives the production cost for a bakery to produce cakes of a certain type, where $C(x)$ is the cost, in dollars, to produce $x$ of the cakes. The function $M$ defined by
  \[
    M(x)=C(x+1)-C(x)
  \]
  gives the marginal cost, in dollars, to produce cake number $x+1$. Which of the following gives the best estimate for the marginal cost, in dollars, to produce the 40th cake?
  \begin{subanswer}
    % Student’s answer space
  \end{subanswer}

  \item \textbf{Left Riemann Sum Approximation} (10 points)\\
  Which of the following is a left Riemann sum approximation of
  \[
    \int_1^7\bigl(4 \ln x + 2\bigr)\,dx
  \]
  with $n$ subintervals of equal length?
  \begin{enumerate}[label=(\Alph*)]
    \item $\displaystyle \sum_{k=1}^n\Bigl(4 \ln\bigl(1+\tfrac{k-1}{n}\bigr)+2\Bigr)\frac{1}{n}$
    \item $\displaystyle \sum_{k=1}^n\Bigl(4 \ln\bigl(\tfrac{6k}{n}\bigr)+2\Bigr)\frac{6}{n}$
    \item $\displaystyle \sum_{k=1}^n\Bigl(4 \ln\bigl(1+\tfrac{6(k-1)}{n}\bigr)+2\Bigr)\frac{6}{n}$
    \item $\displaystyle \sum_{k=1}^n\Bigl(4 \ln\bigl(1+\tfrac{6k}{n}\bigr)+2\Bigr)\frac{6}{n}$
  \end{enumerate}
  \begin{subanswer}
    % Student’s answer space
  \end{subanswer}

  \item \textbf{Instantaneous Rate of Change of the Derivative} (10 points)\\
  Let $f(x)=-\frac{32}{x}+x^3$ on the closed interval $[1,4]$. What is the instantaneous rate of change of $f^{\prime}(x)$ at $x=2$?
  \begin{subanswer}
    % Student’s answer space
  \end{subanswer}

  \item \textbf{Arc Length} (10 points)\\
  What is the length of the curve
  \[
    y = 1 - \cos x
  \]
  from $x=0$ to $x=4\pi$?
  \begin{subanswer}
    % Student’s answer space
  \end{subanswer}

  \item \textbf{Squeeze Theorem Limit} (10 points)\\
  Let $g(x)=-x^2-2x+3$ and $h(x)=\tfrac{1}{2}x^2 + x + \tfrac{13}{2}$. If $f$ is a function satisfying
  \[
    g(x) \le f(x) \le h(x)
    \quad\text{for all }x,
  \]
  what is
  \[
    \lim_{x\to -1} f(x)\;?
  \]
  \begin{subanswer}
    % Student’s answer space
  \end{subanswer}

  \item \textbf{Relative Extrema of an Exponential Polynomial} (10 points)\\
  The function $f$ is defined by
  \[
    f(x)=e^{-x}\bigl(x^2+2x\bigr).
  \]
  At what values of $x$ does $f$ have a relative maximum?
  \begin{subanswer}
    % Student’s answer space
  \end{subanswer}

  \item \textbf{Value from a Given Derivative} (10 points)\\
  Let $f$ be a differentiable function such that $f(0)=5.420$ and
  \[
    f^{\prime}(x)=\sqrt{\sin^2 x + x}.
  \]
  What is the value of $f(2\pi)$?
  \begin{subanswer}
    % Student’s answer space
  \end{subanswer}

  \item \textbf{Limit of a Piecewise Function} (10 points)\\
  Define
  \[
    f(x) =
    \begin{cases}
      \displaystyle \frac{(x-1)^2 (x+1)}{|x-1|}, & x \neq 1,\\[8pt]
      2, & x=1.
    \end{cases}
  \]
  Find
  \[
    \lim_{x\to 1} f(x).
  \]
  \begin{subanswer}
    % Student’s answer space
  \end{subanswer}


  %% continues on ... Q40
  %% continues on ... Q--
  %% continues on ... Q--
  %% continues on ... Q--
  %% continues on ... Q--
  %% continues on ... Q--
  %% continues on ... Q--



  \newpage

  \item \textbf{Indefinite Integrals} (10 points)\\
  Evaluate each of the following indefinite integrals.
  \begin{enumerate}[label=(\alph*)]
    \item $\displaystyle \int \bigl(6 x^5 - 18 x^2 + 7\bigr) \,dx$
    \item $\displaystyle \int 6 x^5 \,dx - 18 x^2 + 7$
  \end{enumerate}
  \begin{subanswer}
    % your answer here
  \end{subanswer}

  \item \textbf{Indefinite Integrals} (10 points)\\
  Evaluate each of the following indefinite integrals.
  \begin{enumerate}[label=(\alph*)]
    \item $\displaystyle \int \bigl(40 x^3 + 12 x^2 - 9 x + 14\bigr) \,dx$
    \item $\displaystyle \int \bigl(40 x^3 + 12 x^2 - 9 x\bigr) \,dx + 14$
    \item $\displaystyle \int \bigl(40 x^3 + 12 x^2\bigr) \,dx - 9 x + 14$
  \end{enumerate}
  \begin{subanswer}
    % your answer here
  \end{subanswer}

  \item \textbf{Indefinite Integral} (5 points)\\
  Evaluate the indefinite integral.
  \[
    \int \bigl(12 t^7 - t^2 - t + 3\bigr) \,dt
  \]
  \begin{subanswer}
    % your answer here
  \end{subanswer}

  \item \textbf{Indefinite Integral} (5 points)\\
  Evaluate the indefinite integral.
  \[
    \int \bigl(10 w^4 + 9 w^3 + 7 w\bigr) \,dw
  \]
  \begin{subanswer}
    % your answer here
  \end{subanswer}

  \item \textbf{Indefinite Integral} (5 points)\\
  Evaluate the indefinite integral.
  \[
    \int \bigl(z^6 + 4 z^4 - z^2\bigr) \,dz
  \]
  \begin{subanswer}
    % your answer here
  \end{subanswer}

  \item \textbf{Function from Derivative} (10 points)\\
  Determine $f(x)$ given that
  \[
    f^{\prime}(x) = 6 x^8 - 20 x^4 + x^2 + 9.
  \]
  \begin{subanswer}
    % your answer here
  \end{subanswer}

  \item \textbf{Function from Derivative} (10 points)\\
  Determine $h(t)$ given that
  \[
    h^{\prime}(t) = t^4 - t^3 + t^2 + t - 1.
  \]
  \begin{subanswer}
    % your answer here
  \end{subanswer}




  \newpage

  \item \textbf{Indefinite Integrals} (105 points)\\
  For problems 1--21, evaluate the given integral.
  \begin{enumerate}
    \item $\displaystyle \int \bigl(4 x^6 - 2 x^3 + 7 x - 4\bigr) \,dx$
    \item $\displaystyle \int \bigl(z^7 - 48 z^{11} - 5 z^{16}\bigr) \,dz$
    \item $\displaystyle \int \bigl(10 t^{-3} + 12 t^{-9} + 4 t^3\bigr) \,dt$
    \item $\displaystyle \int \bigl(w^{-2} + 10 w^{-5} - 8\bigr) \,dw$
    \item $\displaystyle \int 12 \,dy$
    \item $\displaystyle \int \bigl(\sqrt[3]{w} + 10 \sqrt[5]{w^3}\bigr) \,dw$
    \item $\displaystyle \int \bigl(\sqrt{x^7} - 7 \sqrt[6]{x^5} + 17 \sqrt[3]{x^{10}}\bigr) \,dx$
    \item $\displaystyle \int \biggl(\frac{4}{x^2} + 2 - \frac{1}{8 x^3}\biggr) \,dx$
    \item $\displaystyle \int \biggl(\frac{7}{3 y^6} + \frac{1}{y^{10}} - \frac{2}{\sqrt[3]{y^4}}\biggr) \,dy$
    \item $\displaystyle \int \bigl(t^2 - 1\bigr)(4 + 3 t) \,dt$
    \item $\displaystyle \int \sqrt{z}\biggl(z^2 - \frac{1}{4 z}\biggr) \,dz$
    \item $\displaystyle \int \frac{z^8 - 6 z^5 + 4 z^3 - 2}{z^4} \,dz$
    \item $\displaystyle \int \frac{x^4 - \sqrt[3]{x}}{6 \sqrt{x}} \,dx$
    \item $\displaystyle \int \bigl(\sin x + 10 \csc^2 x\bigr) \,dx$
    \item $\displaystyle \int \bigl(2 \cos w - \sec w \tan w\bigr) \,dw$
    \item $\displaystyle \int \bigl(12 + \csc \theta\bigl[\sin \theta + \csc \theta\bigr]\bigr) \,d\theta$
    \item $\displaystyle \int \biggl(4 e^z + 15 - \frac{1}{6 z}\biggr) \,dz$
    \item $\displaystyle \int \biggl(t^3 - \frac{e^{-t} - 4}{e^{-t}}\biggr) \,dt$
    \item $\displaystyle \int \biggl(\frac{6}{w^3} - \frac{2}{w}\biggr) \,dw$
    \item $\displaystyle \int \biggl(\frac{1}{1 + x^2} + \frac{12}{\sqrt{1 - x^2}}\biggr) \,dx$
    \item $\displaystyle \int \biggl(6 \cos z + \frac{4}{\sqrt{1 - z^2}}\biggr) \,dz$
  \end{enumerate}

  \item \textbf{Substitution Method Integrals} (80 points)\\
  For problems 1--16, evaluate the given integral using the substitution method.
  \begin{enumerate}
    \item $\displaystyle \int (8 x - 12)\bigl(4 x^2 - 12 x\bigr)^4 \,dx$
    \item $\displaystyle \int 3 t^{-4}\bigl(2 + 4 t^{-3}\bigr)^{-7} \,dt$
    \item $\displaystyle \int (3 - 4 w)\bigl(4 w^2 - 6 w + 7\bigr)^{10} \,dw$
    \item $\displaystyle \int 5(z - 4) \sqrt[3]{z^2 - 8 z} \,dz$
    \item $\displaystyle \int 90 x^2 \sin\bigl(2 + 6 x^3\bigr) \,dx$
    \item $\displaystyle \int \sec(1 - z) \tan(1 - z) \,dz$
    \item $\displaystyle \int \bigl(15 t^{-2} - 5 t\bigr) \cos\bigl(6 t^{-1} + t^2\bigr) \,dt$
    \item $\displaystyle \int \bigl(7 y - 2 y^3\bigr) e^{y^4 - 7 y^2} \,dy$
    \item $\displaystyle \int \frac{4 w + 3}{4 w^2 + 6 w - 1} \,dw$
    \item $\displaystyle \int \bigl(\cos(3 t) - t^2\bigr)\bigl(\sin(3 t) - t^3\bigr)^5 \,dt$
    \item $\displaystyle \int 4\biggl(\frac{1}{z} - e^{-z}\biggr) \cos\bigl(e^{-z} + \ln z\bigr) \,dz$
    \item $\displaystyle \int \sec^2 v \cdot e^{1 + \tan v} \,dv$
    \item $\displaystyle \int 10 \sin(2 x) \cos(2 x) \sqrt{\cos^2(2 x) + 5} \,dx$
    \item $\displaystyle \int \frac{\csc x \cot x}{2 - \csc x} \,dx$
    \item $\displaystyle \int \frac{6}{7 + y^2} \,dy$
    \item $\displaystyle \int \frac{1}{\sqrt{4 - 9 w^2}} \,dw$
  \end{enumerate}

  \item \textbf{Function from Derivative with Initial Condition} (10 points)\\
  Determine $f(x)$ given that
  \[
    f^{\prime}(x) = 12 x^2 - 4 x
  \]
  and $f(-3) = 17$.
  \begin{subanswer}
    % your answer here
  \end{subanswer}

  \item \textbf{Function from Derivative with Initial Condition} (10 points)\\
  Determine $g(z)$ given that
  \[
    g^{\prime}(z) = 3 z^3 + \frac{7}{2 \sqrt{z}} - e^z
  \]
  and $g(1) = 15 - e$.
  \begin{subanswer}
    % your answer here
  \end{subanswer}

  \item \textbf{Function from Second Derivative with Initial Conditions} (10 points)\\
  Determine $h(t)$ given that
  \[
    h^{\prime \prime}(t) = 24 t^2 - 48 t + 2,
  \]
  $h(1) = -9$, and $h(-2) = -4$.
  \begin{subanswer}
    % your answer here
  \end{subanswer}

  \item \textbf{Substitution Method Integrals} (15 points)\\
  Evaluate each of the following integrals.
  \begin{enumerate}[label=(\alph*)]
    \item $\displaystyle \int \frac{3 x}{1 + 9 x^2} \,dx$
    \item $\displaystyle \int \frac{3 x}{\bigl(1 + 9 x^2\bigr)^4} \,dx$
    \item $\displaystyle \int \frac{3}{1 + 9 x^2} \,dx$
  \end{enumerate}
  \begin{subanswer}
    % your answer here
  \end{subanswer}

  \newpage

  \item \textbf{Substitution Method Integrals} (65 points)\\
  For problems 1--13, evaluate the given integral using the substitution method.
  \begin{enumerate}
    \item $\displaystyle \int \bigl(4 \sqrt{5 + 9 t} + 12(5 + 9 t)^7\bigr) \,dt$
    \item $\displaystyle \int \bigl(7 x^3 \cos\bigl(2 + x^4\bigr) - 8 x^3 e^{2 + x^4}\bigr) \,dx$
    \item $\displaystyle \int \biggl(\frac{6 e^{7 w}}{\bigl(1 - 8 e^{7 w}\bigr)^3} + \frac{14 e^{7 w}}{1 - 8 e^{7 w}}\biggr) \,dw$
    \item $\displaystyle \int \bigl(x^4 - 7 x^5 \cos\bigl(2 x^6 + 3\bigr)\bigr) \,dx$
    \item $\displaystyle \int \biggl(e^z + \frac{4 \sin(8 z)}{1 + 9 \cos(8 z)}\biggr) \,dz$
    \item $\displaystyle \int \bigl(20 e^{2 - 8 w} \sqrt{1 + e^{2 - 8 w}} + 7 w^3 - 6 \sqrt[3]{w}\bigr) \,dw$
    \item $\displaystyle \int \bigl((4 + 7 t)^3 - 9 t \sqrt[4]{5 t^2 + 3}\bigr) \,dt$
    \item $\displaystyle \int \biggl(\frac{6 x - x^2}{x^3 - 9 x^2 + 8} - \csc^2\biggl(\frac{3 x}{2}\biggr)\biggr) \,dx$
    \item $\displaystyle \int \bigl(7(3 y + 2)\bigl(4 y + 3 y^2\bigr)^3 + \sin(3 + 8 y)\bigr) \,dy$
    \item $\displaystyle \int \sec^2(2 t)\bigl[9 + 7 \tan(2 t) - \tan^2(2 t)\bigr] \,dt$
    \item $\displaystyle \int \frac{8 - w}{4 w^2 + 9} \,dw$
    \item $\displaystyle \int \frac{7 x + 2}{\sqrt{1 - 25 x^2}} \,dx$
    \item $\displaystyle \int z^7\bigl(8 + 3 z^4\bigr)^8 \,dz$
  \end{enumerate}
  \begin{subanswer}
    % your answer here
  \end{subanswer}

  \item \textbf{Area Estimation Using Riemann Sums} (45 points)\\
  For problems 1--3, estimate the area of the region between the function and the $x$-axis on the given interval using $n = 6$ and using
  \begin{enumerate}[label=(\alph*)]
    \item the right end points of the subintervals for the height of the rectangles,
    \item the left end points of the subintervals for the height of the rectangles, and
    \item the midpoints of the subintervals for the height of the rectangles.
  \end{enumerate}
  \begin{enumerate}
    \item $f(x) = x^3 - 2 x^2 + 4$ on $[1, 4]$
    \item $g(x) = 4 - \sqrt{x^2 + 2}$ on $[-1, 3]$
    \item $h(x) = -x \cos\bigl(\frac{x}{3}\bigr)$ on $[0, 3]$
  \end{enumerate}
  \begin{subanswer}
    % your answer here
  \end{subanswer}

  \item \textbf{Net Area Estimation} (5 points)\\
  Estimate the net area between $f(x) = 8 x^2 - x^5 - 12$ and the $x$-axis on $[-2, 2]$ using $n = 8$ and the midpoints of the subintervals for the height of the rectangles. Without looking at a graph of the function on the interval, does it appear that more of the area is above or below the $x$-axis?
  \begin{subanswer}
    % your answer here
  \end{subanswer}

  \item \textbf{Definite Integral by Definition} (15 points)\\
  For problems 1--2, use the definition of the definite integral to evaluate the integral. Use the right end point of each interval for $x_i^*$.
  \begin{enumerate}
    \item $\displaystyle \int_1^4 (2 x + 3) \,dx$
    \item $\displaystyle \int_0^1 6 x(x - 1) \,dx$
  \end{enumerate}
  \begin{subanswer}
    % your answer here
  \end{subanswer}

  \item \textbf{Definite Integral Evaluation} (5 points)\\
  Evaluate:
  \[
    \int_4^4 \frac{\cos\bigl(e^{3 x} + x^2\bigr)}{x^4 + 1} \,dx
  \]
  \begin{subanswer}
    % your answer here
  \end{subanswer}

  \item \textbf{Properties of Definite Integrals} (10 points)\\
  For problems 1--2, determine the value of the given integral given that $\int_6^{11} f(x) \,dx = -7$ and $\int_6^{11} g(x) \,dx = 24$.
  \begin{enumerate}
    \item $\displaystyle \int_{11}^6 9 f(x) \,dx$
    \item $\displaystyle \int_6^{11} \bigl(6 g(x) - 10 f(x)\bigr) \,dx$
  \end{enumerate}
  \begin{subanswer}
    % your answer here
  \end{subanswer}

  \item \textbf{Properties of Definite Integrals} (5 points)\\
  Determine the value of $\int_2^9 f(x) \,dx$ given that $\int_5^2 f(x) \,dx = 3$ and $\int_5^9 f(x) \,dx = 8$.
  \begin{subanswer}
    % your answer here
  \end{subanswer}

  \item \textbf{Properties of Definite Integrals} (5 points)\\
  Determine the value of $\int_{-4}^{20} f(x) \,dx$ given that $\int_{-4}^0 f(x) \,dx = -2$, $\int_{31}^0 f(x) \,dx = 19$, and $\int_{20}^{31} f(x) \,dx = -21$.
  \begin{subanswer}
    % your answer here
  \end{subanswer}

  \item \textbf{Area Interpretation of Definite Integrals} (10 points)\\
  For problems 1--2, sketch the graph of the integrand and use the area interpretation of the definite integral to determine the value of the integral.
  \begin{enumerate}
    \item $\displaystyle \int_1^4 (3 x - 2) \,dx$
    \item $\displaystyle \int_0^5 (-4 x) \,dx$
  \end{enumerate}
  \begin{subanswer}
    % your answer here
  \end{subanswer}

  \item \textbf{Fundamental Theorem of Calculus} (15 points)\\
  For problems 1--3, differentiate each of the following integrals with respect to $x$.
  \begin{enumerate}
    \item $\displaystyle \int_4^x 9 \cos^2\bigl(t^2 - 6 t + 1\bigr) \,dt$
    \item $\displaystyle \int_7^{\sin(6 x)} \sqrt{t^2 + 4} \,dt$
    \item $\displaystyle \int_{3 x^2}^{-1} \frac{e^t - 1}{t} \,dt$
  \end{enumerate}
  \begin{subanswer}
    % your answer here
  \end{subanswer}
  \item \textbf{Definite Integrals} (15 points)\\
  Evaluate each of the following integrals.
  \begin{enumerate}[label=(\alph*)]
    \item $\displaystyle \int \biggl(\cos x - \frac{3}{x^5}\biggr) \,dx$
    \item $\displaystyle \int_{-3}^4 \biggl(\cos x - \frac{3}{x^5}\biggr) \,dx$
    \item $\displaystyle \int_1^4 \biggl(\cos x - \frac{3}{x^5}\biggr) \,dx$
  \end{enumerate}
  \begin{subanswer}
    % your answer here
  \end{subanswer}



  \newpage
  \item \textbf{Definite Integrals} (85 points)\\
  Evaluate each of the following integrals, if possible. If it is not possible, clearly explain why it is not possible to evaluate the integral.
  \begin{enumerate}
    \item $\displaystyle \int_1^6 \bigl(12 x^3 - 9 x^2 + 2\bigr) \,dx$
    \item $\displaystyle \int_{-2}^1 \bigl(5 z^2 - 7 z + 3\bigr) \,dz$
    \item $\displaystyle \int_3^0 \bigl(15 w^4 - 13 w^2 + w\bigr) \,dw$
    \item $\displaystyle \int_1^4 \biggl(\frac{8}{\sqrt{t}} - 12 \sqrt{t^3}\biggr) \,dt$
    \item $\displaystyle \int_1^2 \biggl(\frac{1}{7 z} + \frac{\sqrt[3]{z^2}}{4} - \frac{1}{2 z^3}\biggr) \,dz$
    \item $\displaystyle \int_{-2}^4 \biggl(x^6 - x^4 + \frac{1}{x^2}\biggr) \,dx$
    \item $\displaystyle \int_{-4}^{-1} x^2(3 - 4 x) \,dx$
    \item $\displaystyle \int_2^1 \frac{2 y^3 - 6 y^2}{y^2} \,dy$
    \item $\displaystyle \int_0^{\frac{\pi}{2}} \bigl(7 \sin t - 2 \cos t\bigr) \,dt$
    \item $\displaystyle \int_0^\pi \bigl(\sec z \tan z - 1\bigr) \,dz$
    \item $\displaystyle \int_{\frac{\pi}{6}}^{\frac{\pi}{3}} \bigl(2 \sec^2 w - 8 \csc w \cot w\bigr) \,dw$
    \item $\displaystyle \int_0^2 \biggl(e^x + \frac{1}{x^2 + 1}\biggr) \,dx$
    \item $\displaystyle \int_{-5}^{-2} \biggl(7 e^y + \frac{2}{y}\biggr) \,dy$
    \item $\displaystyle \int_0^4 f(t) \,dt$ where
      \[
        f(t) = \begin{cases}
          2 t & t > 1, \\
          1 - 3 t^2 & t \leq 1.
        \end{cases}
      \]
    \item $\displaystyle \int_{-6}^1 g(z) \,dz$ where
      \[
        g(z) = \begin{cases}
          2 - z & z > -2, \\
          4 e^z & z \leq -2.
        \end{cases}
      \]
    \item $\displaystyle \int_3^6 |2 x - 10| \,dx$
    \item $\displaystyle \int_{-1}^0 |4 w + 3| \,dw$
  \end{enumerate}

  \item \textbf{Definite Integrals Using Substitution} (50 points)\\
  Evaluate each of the following integrals, if possible. If it is not possible, clearly explain why it is not possible to evaluate the integral.
  \begin{enumerate}
    \item $\displaystyle \int_0^1 3\bigl(4 x + x^4\bigr)\bigl(10 x^2 + x^5 - 2\bigr)^6 \,dx$
    \item $\displaystyle \int_0^{\frac{\pi}{4}} \frac{8 \cos(2 t)}{\sqrt{9 - 5 \sin(2 t)}} \,dt$
    \item $\displaystyle \int_\pi^0 \sin z \cos^3 z \,dz$
    \item $\displaystyle \int_1^4 \sqrt{w} e^{1 - \sqrt{w^3}} \,dw$
    \item $\displaystyle \int_{-4}^{-1} \biggl(\sqrt[3]{5 - 2 y} + \frac{7}{5 - 2 y}\biggr) \,dy$
    \item $\displaystyle \int_{-1}^2 \biggl(x^3 + e^{\frac{1}{4} x}\biggr) \,dx$
    \item $\displaystyle \int_\pi^{\frac{3 \pi}{2}} \bigl(6 \sin(2 w) - 7 \cos w\bigr) \,dw$
    \item $\displaystyle \int_1^5 \biggl(\frac{2 x^3 + x}{x^4 + x^2 + 1} - \frac{x}{x^2 - 4}\biggr) \,dx$
    \item $\displaystyle \int_{-2}^0 \biggl(t \sqrt{3 + t^2} + \frac{3}{(6 t - 1)^2}\biggr) \,dt$
    \item $\displaystyle \int_{-2}^1 \bigl((2 - z)^3 + \sin(\pi z)\bigl[3 + 2 \cos(\pi z)\bigr]^3\bigr) \,dz$
  \end{enumerate}
  \begin{subanswer}
    % your answer here
  \end{subanswer}


  % https://tutorial.math.lamar.edu/Problems/CalcI/IntAppsIntro.aspx

\end{enumerate}
